\section*{الفصل \ref{ch:b2:metamorphic-facies} --- التفاعلات التحوّلية والسحنات}

\begin{solution}{exo:b2:metamorphic-facies:1}
(أ) تحوّل تماسي، وهو تحوّل حراري محلّي بحرارة المتداخل.
(ب) تحوّل الانزلاق، في نطاق صدع: تصغير حجم الحبيبات، وميلونيت.
(ج) تحوّل قاع المحيط، بمياه البحر الساخنة قرب الظهرة، مع إعادة تبلور
غير مكتملة. (د) تحوّل الدفن، وهو إقليمي، دون تشوّه ولا
شستية.
\end{solution}

\begin{solution}{exo:b2:metamorphic-facies:2}
$\Delta_r S = 95.40 - 82.80 = \qty{12.60}{J.K^{-1}}$، $\Delta_r V = 4.986 - 4.415 =
\qty{0.571}{J.bar^{-1}}$؛ والميل $12.60/0.571 = \qty{22.1}{bar/K} = \qty{2.21}{MPa/K}$.
وعلى مدى \qty{100}{\celsius} يرتفع ضغط التوازن بمقدار \qty{2.21}{kbar}.
\end{solution}

\begin{solution}{exo:b2:metamorphic-facies:3}
(أ) الشيست الأخضر. (ب) الشيست الأزرق. (ج) الأمفيبوليت: هورنبلند مع بلاجيوكليز فيه أكثر
من \qty{17}{\percent} أنورثيت. (د) الإكلوجيت.
\end{solution}

\begin{solution}{exo:b2:metamorphic-facies:4}
(أ) $\rho = M/V$: الكيانيت $162.046/44.15 = \qty{3.670}{g.cm^{-3}}$، والأندالوسيت
$162.046/51.52 = 3.145$، والسيليمانيت $162.046/49.86 = 3.250$. (ب) الكيانيت: عند ضغط
أعلى ينمو الحدّ $V\,(P - 1)$ من طاقة جيبس أقلّ ما ينمو لأصغر حجم.
(ج) تحوي البلّورات الطبيعية عناصر أخرى (فالحديد الثلاثي، مثلًا، يحلّ محلّ بعض
الألومنيوم في الأندالوسيت والسيليمانيت الطبيعيين)، والقياسات على العيّنات
الطبيعية أقل دقة من الأحجام المولية في الجداول.
\end{solution}

\begin{solution}{exo:b2:metamorphic-facies:5}
(أ) $\Delta_r G = 3900 - 298.15\times8.59 = \qty{+1339}{J} > 0$: الكيانيت. (ب) عند
\qty{1}{bar}، $\Delta_r H = T\Delta_r S$: $T = 3900/8.59 = \qty{454}{K} =
\qty{181}{\celsius}$. (ج) $P = (8.59\times673.15 - 3900)/0.737 = 1882/0.737 =
\qty{2554}{bar} = \qty{2.55}{kbar}$، مقابل 2.54 مقيسة. (د) التفاعل أبطأ من أن
يجري عند درجة حرارة منخفضة: فيبقى الأندالوسيت شبه مستقر، كما تبقى معظم معادن درجات الحرارة
العالية على السطح.
\end{solution}

\begin{solution}{exo:b2:metamorphic-facies:6}
(أ) $\Delta_r H = 7700 + 1000\times(121.39 - 122.29) = \qty{6800}{J}$؛ $\Delta_r S = 296.44
- 285.02 = \qty{11.42}{J.K^{-1}}$؛ $\Delta_r G = 6800 - 1000\times11.42 = \qty{-4620}{J}$.
(ب) بالقيم عند \qty{298.15}{K}: $7700 - 1000\times12.60 = \qty{-4900}{J}$.
(ج) السيليمانيت. تخطئ الطريقة البسيطة بمقدار \qty{280}{J}، أي نحو \qty{6}{\percent}؛ وبالقسمة
على $\Delta_r V = 0.571$، يكون ذلك نحو \qty{490}{bar} على الخط، وهو أقل بكثير من
أثر عدم اليقين في الإنثالبيات.
\end{solution}

\begin{solution}{exo:b2:metamorphic-facies:7}
(أ) $\Delta_r G = -1127.4 + 1128.5 = \qty{1.1}{kJ}$ و$\Delta_r V = (34.15 -
36.93)/10 = \qty{-0.278}{J.bar^{-1}}$: $P = 1100/0.278 = \qty{3957}{bar} \approx
\qty{4.0}{kbar}$. (ب) $\Delta_r S = 88.0 - 91.7 = -3.7$؛ والميل $-3.7/-0.278 =
\qty{13.3}{bar/K}$. (ج) مع $\Delta_r H = 0$، $P = T\Delta_r S/\Delta_r V = 13.31\,T$:
عند \qty{573.15}{K}، \qty{7629}{bar} $= \qty{7.63}{kbar}$. (د) أراغونيت الأصداف
شبه مستقر على السطح: تصنعه الكائنات الحية، ويبقى لأنه لا يتحوّل إلى
كالسيت إلا ببطء.
\end{solution}

\begin{solution}{exo:b2:metamorphic-facies:8}
(أ) $-32.4/-1.698 = \qty{19.1}{bar/K}$. (ب) $P = (T\Delta_r S - \Delta_r H)/\Delta_r V =
(-32.4\times773.15 + 5000)/(-1.698) = (-25050 + 5000)/(-1.698) = \qty{11808}{bar}
\approx \qty{11.8}{kbar}$. (ج) $\sqrt{2.6^2 + 3.6^2 + 1.0^2} = \qty{4.55}{kJ}$ على
$\Delta_r H$، ومنه $4550/1.698 = \qty{2.7}{kbar}$ على الضغط. (د) عند ضغوط
سحنة الإكلوجيت يكون مكوّن الألبيت في البلاجيوكليز غير مستقر بالنسبة إلى
الجاديت والكوارتز؛ فيدخل الصوديوم في البيروكسين (الأومفاسيت)، ويختفي
البلاجيوكليز.
\end{solution}

\begin{solution}{exo:b2:metamorphic-facies:9}
(أ) من 1 إلى \qty{2}{kbar}: $1000/34 = \qty{29.4}{bar/K}$؛ ومن 2 إلى 3: $1000/25 =
\qty{40.0}{bar/K}$. (ب) مع ثبات $\Delta_r S$، يتناسب $\Delta_r V$ عكسيًا مع
الميل: فقد تقلّص بعامل $40.0/29.4 = 1.36$. (ج) الماء المتحرّر: إذ ينخفض حجمه
المولي كلما ارتفع الضغط، بينما لا تكاد أحجام المواد الصلبة تتغيّر.
\end{solution}

\begin{solution}{exo:b2:metamorphic-facies:10}
(أ) عند \qty{873.15}{K}: الكيانيت $=$ الأندالوسيت $(8.59\times873.15 - 3900)/0.737 =
\qty{4885}{bar}$؛ والكيانيت $=$ السيليمانيت $(12.60\times873.15 - 7700)/0.571 =
\qty{5782}{bar}$؛ والأندالوسيت $=$ السيليمانيت $(4.01\times873.15 - 3800)/(-0.166) =
\qty{1799}{bar}$. (ب) عند \qty{4885}{bar}، تحت خط الكيانيت $=$ السيليمانيت (5782)،
يكون السيليمانيت أكثر استقرارًا من الكيانيت، وفوق خط الأندالوسيت $=$ السيليمانيت
(1799)، أكثر استقرارًا من الأندالوسيت: فالنقطة تقع في حقل السيليمانيت،
وتوازن الكيانيت $=$ الأندالوسيت هناك شبه مستقر. (ج) عند \qty{4}{kbar}: بين 1799
و\qty{5782}{bar}، السيليمانيت. وعند \qty{7}{kbar}: فوق \qty{5782}{bar} يتغلّب الكيانيت على
السيليمانيت، وفوق \qty{4885}{bar} يتغلّب على الأندالوسيت: الكيانيت.
\end{solution}

\begin{solution}{exo:b2:metamorphic-facies:11}
(أ) $z = P/(\rho g) = 3.75\times10^8/(2700\times9.81) = \qty{14.2}{km}$، بعدم يقين قدره
$0.25\times10^8/26487 = \qty{0.94}{km}$. (ب) $(504 - 10)/14.16 =
\qty{35}{\celsius.km^{-1}}$. (ج) لا تتعايش المعادن الثلاثة في توازن إلا عند
\omterm{def:b2:metamorphic-facies:invariant}{النقطة اللامتغيّرة}: فلا بدّ أن يمرّ مسار الصخر بها، وأن تتوقّف التفاعلات
عندها. ومعظم الصخور تمرّ بجانبها وتحتفظ بمتعدّد شكل واحد أو اثنين.
\end{solution}

\begin{solution}{exo:b2:metamorphic-facies:12}
(أ) X: $T = 10 + 10\times30 = \qty{310}{\celsius}$ (\qty{583.15}{K})؛ $P = 3000\times
9.81\times30000 = \qty{8.83e8}{Pa} = \qty{8.83}{kbar}$. Y: $T = 10 + 30\times20 =
\qty{610}{\celsius}$ (\qty{883.15}{K})؛ $P = 2800\times9.81\times20000 = \qty{5.49}{kbar}$.
(ب) X: الكالسيت $=$ الأراغونيت عند $13.31\times583.15 = \qty{7.76}{kbar} < 8.83$، فالأراغونيت؛
والألبيت $=$ الجاديت + الكوارتز عند $(32.4\times583.15 - 5000)/1.698 = \qty{8.18}{kbar} <
8.83$، فالجاديت + الكوارتز؛ والكيانيت $=$ الأندالوسيت عند $(8.59\times583.15 - 3900)/0.737 =
\qty{1.51}{kbar}$، تحت درجة حرارة النقطة الثلاثية: فالكيانيت. Y: الكالسيت $=$
الأراغونيت عند \qty{11.75}{kbar} $> 5.49$، فالكالسيت؛ والألبيت $=$ الجاديت + الكوارتز عند
\qty{13.9}{kbar}، فالألبيت؛ وفوق النقطة الثلاثية، الكيانيت $=$ السيليمانيت عند
\qty{6.00}{kbar} $> 5.49$ والأندالوسيت $=$ السيليمانيت عند \qty{1.56}{kbar} $< 5.49$:
فالسيليمانيت. (ج) المسار X، بمعدل \qty{10}{\celsius\per\km}، هو الأبرد: فالجاديت
والأراغونيت عند \qty{310}{\celsius} فقط يضعانه في سحنة الشيست الأزرق. أمّا Y، الأدفأ بثلاث
مرات، فينتهي بالسيليمانيت عند \qty{610}{\celsius} وضغط معتدل، في
سحنة الأمفيبوليت.
\end{solution}

\begin{solution}{pb:b2:metamorphic-facies:1}
\noindent\textbf{الجزء الأول.}

\textbf{1.} $\Delta_r H = -2589.9 + 2593.8 = \qty{3.9}{kJ}$؛ $\Delta_r S = 91.39 - 82.80 =
\qty{8.59}{J.K^{-1}}$؛ $\Delta_r V = 5.152 - 4.415 = \qty{0.737}{J.bar^{-1}}$.

\textbf{2.} $\qty{7.7}{kJ}$؛ $\qty{12.60}{J.K^{-1}}$؛ $\qty{0.571}{J.bar^{-1}}$.

\textbf{3.} $\qty{3.8}{kJ}$؛ $\qty{4.01}{J.K^{-1}}$؛ $\qty{-0.166}{J.bar^{-1}}$؛ وكلٌّ منها
هو الثاني ناقص الأول، لأن الأندالوسيت يُختزل.

\textbf{4.} للسيليمانيت أكبر إنتروبيا (95.40): فهو يقع في جهة درجة الحرارة المرتفعة
من خطّيه كليهما. وللكيانيت أصغر حجم (4.415): فهو يقع في جهة
الضغط المرتفع من خطّيه كليهما. والأندالوسيت، صاحب أكبر حجم، يشغل
جهة الضغط المنخفض.

\textbf{5.} $8.59/0.737 = \qty{11.66}{bar/K}$؛ $12.60/0.571 = \qty{22.07}{bar/K}$؛
$4.01/(-0.166) = \qty{-24.16}{bar/K}$.

\textbf{6.} للسيليمانيت إنتروبيا أكبر وحجم أصغر من الأندالوسيت: $\Delta_r S > 0$
و$\Delta_r V < 0$. فرفع درجة الحرارة يفضّل السيليمانيت، ورفع الضغط
كذلك، فيتّجه الخط إلى ضغط أدنى عند درجة حرارة أعلى.

\noindent\textbf{الجزء الثاني.}

\textbf{7.} $3900 - 298.15\times8.59 = \qty{+1339}{J}$؛ $7700 - 298.15\times12.60 =
\qty{+3943}{J}$. كلاهما موجب: فالكيانيت مستقر على السطح.

\textbf{8.} $T = \Delta_r H/\Delta_r S$: $3900/8.59 = \qty{454.0}{K} = \qty{181}{\celsius}$؛
$7700/12.60 = \qty{611.1}{K} = \qty{338}{\celsius}$؛ $3800/4.01 = \qty{947.6}{K} =
\qty{674}{\celsius}$.

\textbf{9.} بالتسخين عند \qty{1}{bar}، يُفسح الكيانيت المجال للأندالوسيت عند
\qty{181}{\celsius}؛ وعند \qty{338}{\celsius} لا يعود الكيانيت مستقرًا، فيكون توازن الكيانيت $=$
السيليمانيت شبه مستقر هناك؛ ويُفسح الأندالوسيت المجال للسيليمانيت عند
\qty{674}{\celsius}. فالأول والثالث توازنان مستقران.

\textbf{10.} $(8.59\times673.15 - 3900)/0.737 = \qty{2554}{bar} = \qty{2.55}{kbar}$،
على بُعد \qty{0.01}{kbar} من التجربة.

\textbf{11.} $(12.60\times973.15 - 7700)/0.571 = 4562/0.571 = \qty{7989}{bar} =
\qty{7.99}{kbar}$، أعلى من التجربة بمقدار \qty{0.13}{kbar}.

\textbf{12.} يصعد خط الكيانيت $=$ الأندالوسيت من \qty{181}{\celsius} عند \qty{1}{bar} بميل
\qty{11.7}{bar/K}؛ وينزل خط الأندالوسيت $=$ السيليمانيت من نحو \qty{11.5}{kbar} عند
\qty{200}{\celsius} إلى \qty{1}{bar} عند \qty{674}{\celsius}؛ وخط الكيانيت $=$ السيليمانيت،
الأشدّ انحدارًا بين الخطين الصاعدين، يقطع كليهما قرب \qty{514}{\celsius} و\qty{3.9}{kbar}، وكل خط مستقر في
جهة واحدة فقط من التقاطع (مخطّط الفصل).

\noindent\textbf{الجزء الثالث.}

\textbf{13.} $a_1 = 8.59/0.737 = \qty{11.655}{bar/K}$، $b_1 = 3900/0.737 =
\qty{5291.7}{bar}$؛ $a_2 = 12.60/0.571 = \qty{22.067}{bar/K}$، $b_2 = 7700/0.571 =
\qty{13485.1}{bar}$.

\textbf{14.} $T = (b_2 - b_1)/(a_2 - a_1) = 8193.4/10.412 = \qty{786.92}{K} =
\qty{513.8}{\celsius}$.

\textbf{15.} $P = 11.655\times786.92 - 5291.7 = \qty{3880}{bar} = \qty{3.88}{kbar}$.

\textbf{16.} $(4.01\times786.92 - 3800)/(-0.166) = (3155.5 - 3800)/(-0.166) =
\qty{3882}{bar}$: الضغط نفسه في حدود تقريب البيانات، كما تقتضي قضية الفصل.

\textbf{17.} تقع \qty{514}{\celsius} على بُعد \qty{10}{\celsius} من \qty{504}{\celsius}، داخل
عدم اليقين البالغ \qty{20}{\celsius}؛ وتقع \qty{3.88}{kbar} على بُعد \qty{0.13}{kbar} من 3.75،
داخل 0.25. فالتوافق جيد.

\textbf{18.} لا: تقع \qty{536}{\celsius} على بُعد \qty{32}{\celsius} من التجربة، خارج
عدم يقينها. فإهمال السعات الحرارية لم يكن الخطأ الرئيسي: بل الإنثالبيات.

\noindent\textbf{الجزء الرابع.}

\textbf{19.} ترتفع $\Delta_r H$ للتفاعل كيانيت $\to$ أندالوسيت بمقدار \qty{1000}{J}؛ وعند
$T$ ثابتة، تنخفض $P = (T\Delta_r S - \Delta_r H)/\Delta_r V$ بمقدار $1000/0.737 = \qty{1357}{bar}$.

\textbf{20.} $\Delta T = -1357/(a_2 - a_1) = -1357/10.41 = \qty{-130}{K}$: فتنتقل النقطة الثلاثية
إلى نحو \qty{384}{\celsius}.

\textbf{21.} ينخفض خط الكيانيت $=$ السيليمانيت بمقدار $1000/0.571 = \qty{1751}{bar}$،
وتنتقل النقطة الثلاثية بمقدار $+1751/10.41 = \qty{+168}{K}$.

\textbf{22.} كل إنثالبي وحدها، إذا غُيّرت بمقدار $\pm\qty{2.0}{kJ}$، تنقل النقطة الثلاثية بمقدار
$\pm\qty{260}{K}$ (الأندالوسيت) أو $\pm\qty{336}{K}$ (السيليمانيت): فقياس الحرارة وحده يسمح
بنقاط ثلاثية تمتدّ على عدة مئات من الدرجات.

\textbf{23.} الإنثالبيات المجدولة ليست مستمدّة من قياس الحرارة وحده: ففي حدود
عدم يقينها تُضبط بحيث تلائم الخطوط المحسوبة التجارب المعكوسة.
فالتجارب تثبّت المخطّط؛ والجداول تنقله، على نحو متّسق، إلى تفاعلات أخرى.

\textbf{24.} تقع النقطة الثلاثية المحسوبة من الجداول بخواص ثابتة عند
\textbf{\qty[text-series-to-math=true]{514}{\celsius}}.
\end{solution}
